% ====================================================================
% ФАЙЛ: tasks/01_mechanics/momentum_bank.tex
% ТЕМА: ИМПУЛЬС ТЕЛА. ЗАКОН СОХРАНЕНИЯ ИМПУЛЬСА
% ПАЧКА 1: ВАРИАНТЫ 1–10
% ====================================================================

% === ЗАДАЧА 1 ===
% Тег: #МЕХАНИКА #КИМ_03 #ВАРИАНТ_01
\begin{taskbasic}[code={\label{task:momentum_v1_n3}}]
\textbf{Задача 1.} Тело массой $2$~кг движется вдоль оси $Ox$. Зависимость проекции скорости $v_x$ этого тела от времени $t$ имеет вид: $v_x = 5 + 3t$ (все величины выражены в СИ). Чему равен модуль изменения импульса тела за первые $2$~с движения?

\kim{3}
\end{taskbasic}
\answer{6}

% === ЗАДАЧА 2 ===
% Тег: #МЕХАНИКА #КИМ_03 #ВАРИАНТ_02
\begin{taskbasic}[code={\label{task:momentum_v2_n3}}]
\textbf{Задача 2.} Тело массой $1$~кг движется вдоль оси $Ox$. Зависимость проекции скорости $v_x$ этого тела от времени $t$ имеет вид: $v_x = 2 + 4t$ (все величины выражены в СИ). Чему равен модуль изменения импульса тела за первую $1$~с движения?

\kim{3}
\end{taskbasic}
\answer{4}

% === ЗАДАЧА 3 ===
% Тег: #МЕХАНИКА #КИМ_03 #ВАРИАНТ_05
\begin{taskbasic}[code={\label{task:momentum_v5_n3}}]
\textbf{Задача 3.} В инерциальной системе отсчёта тело движется по прямой линии в одном направлении под действием постоянной силы. На рисунке приведён график зависимости модуля этой силы $F$ от времени $t$. Чему равен модуль изменения импульса тела в интервале времени от $0$ до $3$~с?

\nopagebreak\smallskip
\begin{center}
\begin{tikzpicture}[x=1.5cm, y=0.8cm, every node/.style={font=\footnotesize}]
    \draw[xstep=1, ystep=1, gray!30, thin] (0,0) grid (4,5);
    \draw[-{Stealth[scale=0.8]}, black, thick] (0,0) -- (4.5,0) node[right] {$t, \text{ с}$};
    \draw[-{Stealth[scale=0.8]}, black, thick] (0,0) -- (0,5.5) node[above] {$F, \text{ Н}$};
    \foreach \x in {1,2,3,4} \node[below] at (\x,0) {\x};
    \foreach \y in {1,2,3,4,5} \node[left] at (0,\y) {\y};
    \node[below left] at (0,0) {0};
    \draw[ultra thick, brandPrimary] (0,0) -- (2,4) -- (3,4) -- (4,0);
    \fill[black] (2,4) circle (1.5pt);
    \fill[black] (3,4) circle (1.5pt);
\end{tikzpicture}
\end{center}

\kim{3}
\end{taskbasic}
\answer{8}

% === ЗАДАЧА 4 ===
% Тег: #МЕХАНИКА #КИМ_03 #ВАРИАНТ_06
\begin{taskbasic}[code={\label{task:momentum_v6_n3}}]
\textbf{Задача 4.} В инерциальной системе отсчёта тело движется по прямой линии в одном направлении под действием постоянной силы. На рисунке приведён график зависимости модуля этой силы $F$ от времени $t$. Чему равен модуль изменения импульса тела в интервале времени от $0$ до $4$~с?

\nopagebreak\smallskip
\begin{center}
\begin{tikzpicture}[x=1.5cm, y=0.8cm, every node/.style={font=\footnotesize}]
    \draw[xstep=1, ystep=1, gray!30, thin] (0,0) grid (5,5);
    \draw[-{Stealth[scale=0.8]}, black, thick] (0,0) -- (5.5,0) node[right] {$t, \text{ с}$};
    \draw[-{Stealth[scale=0.8]}, black, thick] (0,0) -- (0,5.5) node[above] {$F, \text{ Н}$};
    \foreach \x in {1,2,3,4,5} \node[below] at (\x,0) {\x};
    \foreach \y in {1,2,3,4,5} \node[left] at (0,\y) {\y};
    \node[below left] at (0,0) {0};
    \draw[ultra thick, brandPrimary] (0,0) -- (2,4) -- (4,4) -- (5,0);
    \fill[black] (2,4) circle (1.5pt);
    \fill[black] (4,4) circle (1.5pt);
\end{tikzpicture}
\end{center}

\kim{3}
\end{taskbasic}
\answer{12}

% === ЗАДАЧА 5 ===
% Тег: #МЕХАНИКА #КИМ_03 #ВАРИАНТ_09
\begin{taskbasic}[code={\label{task:momentum_v9_n3}}]
\textbf{Задача 5.} Железнодорожный вагон массой $20$~т, движущийся со скоростью $0{,}5$~м/с, сталкивается с неподвижной платформой массой $60$~т и сцепляется с ней. Чему равна скорость совместного движения вагона и платформы после сцепки?

\kim{3}
\end{taskbasic}
\answer{0{,}125}

% === ЗАДАЧА 6 ===
% Тег: #МЕХАНИКА #КИМ_03 #ВАРИАНТ_10
\begin{taskbasic}[code={\label{task:momentum_v10_n3}}]
\textbf{Задача 6.} Снаряд массой $2$~кг, летящий со скоростью $20$~м/с, разрывается на два осколка. Один осколок массой $1$~кг летит после взрыва в прежнем направлении со скоростью $25$~м/с. Чему равен модуль скорости второго осколка?

\kim{3}
\end{taskbasic}
\answer{15}

% ====================================================================
% ФАЙЛ: tasks/01_mechanics/momentum_bank_part2.tex
% ТЕМА: ИМПУЛЬС ТЕЛА. ЗАКОН СОХРАНЕНИЯ ИМПУЛЬСА
% ПАЧКА 2: ВАРИАНТЫ 11–20
% ====================================================================

% === ЗАДАЧА 1 ===
% Тег: #МЕХАНИКА #КИМ_03 #ВАРИАНТ_11
\begin{taskbasic}[code={\label{task:momentum_v11_n3}}]
\textbf{Задача 1.} Шарик массой $200$~г падает с высоты $5$~м на горизонтальную плиту. Чему равен модуль изменения импульса шарика при абсолютно упругом ударе? Сопротивлением воздуха пренебречь.

\kim{3}
\end{taskbasic}
\answer{4}

% === ЗАДАЧА 2 ===
% Тег: #МЕХАНИКА #КИМ_03 #ВАРИАНТ_12
\begin{taskbasic}[code={\label{task:momentum_v12_n3}}]
\textbf{Задача 2.} Шарик массой $300$~г падает с высоты $5$~м на горизонтальную плиту. Чему равен модуль изменения импульса шарика при абсолютно упругом ударе? Сопротивлением воздуха пренебречь.

\kim{3}
\end{taskbasic}
\answer{6}

% === ЗАДАЧА 3 ===
% Тег: #МЕХАНИКА #КИМ_03 #ВАРИАНТ_15
\begin{taskbasic}[code={\label{task:momentum_v15_n3}}]
\textbf{Задача 3.} Тело массой $0.1$~кг влетает в зазор между двумя стенками со скоростью $15$~м/с, направленной под углом к нормали одной из стенок, и упруго отскакивает от неё. Чему равен модуль изменения импульса тела при ударе, если модуль его скорости после удара не изменился?

\kim{3}
\end{taskbasic}
\answer{3}

% === ЗАДАЧА 4 ===
% Тег: #МЕХАНИКА #КИМ_03 #ВАРИАНТ_16
\begin{taskbasic}[code={\label{task:momentum_v16_n3}}]
\textbf{Задача 4.} Два тела движутся по взаимно перпендикулярным направлениям. Модуль импульса первого тела равен $12$~кг$\cdot$м/с, а второго — $16$~кг$\cdot$м/с. Чему равен модуль импульса системы этих двух тел после их абсолютно неупругого соударения?

\kim{3}
\end{taskbasic}
\answer{20}

% === ЗАДАЧА 5 ===
% Тег: #МЕХАНИКА #КИМ_06 #ВАРИАНТ_19
\begin{taskbasic}[code={\label{task:momentum_v19_n6}}]
\textbf{Задача 5.} Стальной шарик падает с некоторой высоты на горизонтальную плиту и упруго отскакивает от неё вверх. Как изменяются в результате удара модуль изменения импульса шарика за время удара и его кинетическая энергия сразу после отскока по сравнению с кинетической энергией перед ударом?

Для каждой величины определите соответствующий характер изменения:
1) увеличивается;
2) уменьшается;
3) не изменяется.

Запишите в таблицу выбранные цифры для каждой физической величины. Цифры в ответе могут повторяться.
\begin{center}
\begin{tabular}{|c|c|}
\hline
Модуль изменения импульса & Кинетическая энергия после отскока \\ \hline
 & \\ \hline
\end{tabular}
\end{center}

\kim{6}
\end{taskbasic}
\answer{12}

% === ЗАДАЧА 6 ===
% Тег: #МЕХАНИКА #КИМ_06 #ВАРИАНТ_20
\begin{taskbasic}[code={\label{task:momentum_v20_n6}}]
\textbf{Задача 6.} Пластилиновый шарик падает с некоторой высоты на горизонтальную плиту и прилипает к ней. Как изменяются в результате удара модуль изменения импульса шарика за время удара и его полная механическая энергия сразу после удара?

Для каждой величины определите соответствующий характер изменения:
1) увеличивается;
2) уменьшается;
3) не изменяется.

Запишите в таблицу выбранные цифры для каждой физической величины. Цифры в ответе могут повторяться.
\begin{center}
\begin{tabular}{|c|c|}
\hline
Модуль изменения импульса & Полная механическая энергия \\ \hline
 & \\ \hline
\end{tabular}
\end{center}

\kim{6}
\end{taskbasic}
\answer{21}
% ====================================================================
% ФАЙЛ: tasks/01_mechanics/momentum_bank_part3.tex
% ТЕМА: ИМПУЛЬС ТЕЛА. ЗАКОН СОХРАНЕНИЯ ИМПУЛЬСА
% ПАЧКА 3: ВАРИАНТЫ 21–30
% ====================================================================


% === ЗАДАЧА 3 ===
% Тег: #МЕХАНИКА #КИМ_03 #ВАРИАНТ_27
\begin{taskbasic}[code={\label{task:momentum_v27_n3}}]
\textbf{Задача 3.} Два тела движутся по взаимно перпендикулярным направлениям. Модуль импульса первого тела равен $9$~кг$\cdot$м/с. Каков модуль импульса второго тела, если после их абсолютно неупругого соударения модуль импульса системы этих двух тел стал равен $15$~кг$\cdot$м/с?

\kim{3}
\end{taskbasic}
\answer{12}

% === ЗАДАЧА 4 ===
% Тег: #МЕХАНИКА #КИМ_03 #ВАРИАНТ_28
\begin{taskbasic}[code={\label{task:momentum_v28_n3}}]
\textbf{Задача 4.} Два тела движутся по взаимно перпендикулярным направлениям. Модуль импульса второго тела равен $12$~кг$\cdot$м/с. Каков модуль импульса первого тела, если после их абсолютно неупругого соударения модуль импульса системы этих двух тел стал равен $20$~кг$\cdot$м/с?

\kim{3}
\end{taskbasic}
\answer{16}

% ====================================================================
% ФАЙЛ: tasks/01_mechanics/momentum_hidden_part2_bank.tex
% ТЕМА: ЗАКОН СОХРАНЕНИЯ ИМПУЛЬСА (ПОЛНЫЙ НАБОР ЧАСТИ 2)
% ====================================================================

% === ЗАДАЧА 1 ===
% Тег: #МЕХАНИКА #КИМ_22 #ВАРИАНТ_01
\begin{taskbasic}[code={\label{task:momentum_v1_n22}}]
\textbf{Задача 1.} Столкнулись два одинаковых пластилиновых шарика, движущихся по гладкой горизонтальной поверхности, причём векторы их скоростей непосредственно перед столкновением были взаимно перпендикулярны и втрие различались по модулю: $v_2 = 3v_1$. Какова скорость шариков после абсолютно неупругого столкновения, если перед столкновением скорость более быстрого шарика была равна по модулю $6$~м/с?

\kim{22}
\end{taskbasic}
\answer{3{,}16~м/с}

% === ЗАДАЧА 2 ===
% Тег: #МЕХАНИКА #КИМ_22 #ВАРИАНТ_02
\begin{taskbasic}[code={\label{task:momentum_v2_n22}}]
\textbf{Задача 2.} Столкнулись два одинаковых пластилиновых шарика, движущихся по гладкой горизонтальной поверхности, причём векторы их скоростей непосредственно перед столкновением были взаимно перпендикулярны и вдвое различались по модулю: $v_2 = 2v_1$. Какова скорость более медленного шарика перед столкновением, если после абсолютно неупругого столкновения их скорость стала равна по модулю $4{,}5$~м/с?

\kim{22}
\end{taskbasic}
\answer{4{,}47~м/с}

% === ЗАДАЧА 3 ===
% Тег: #МЕХАНИКА #КИМ_22 #ВАРИАНТ_03
\begin{taskbasic}[code={\label{task:momentum_v3_n22}}]
\textbf{Задача 3.} Два пластилиновых шарика массами $m_1 = 15$~г и $m_2 = 25$~г, летящие навстречу друг другу с одинаковыми по модулю скоростями $v_1 = v_2 = 4$~м/с, при столкновении слипаются. Какой становится их скорость сразу после столкновения? Считать, что время взаимодействия шариков мало.

\kim{22}
\end{taskbasic}
\answer{1~м/с}

% === ЗАДАЧА 4 ===
% Тег: #МЕХАНИКА #КИМ_22 #ВАРИАНТ_04
\begin{taskbasic}[code={\label{task:momentum_v4_n22}}]
\textbf{Задача 4.} Два пластилиновых шарика массами $m_1 = 20$~г и $m_2 = 30$~г, летящие навстречу друг другу со скоростями $v_1 = 4$~м/с и $v_2 = 3$~м/с, при столкновении слипаются. Какой становится их скорость сразу после столкновения? Считать, что время взаимодействия шариков мало.

\kim{22}
\end{taskbasic}
\answer{0{,}2~м/с}

% === ЗАДАЧА 5 ===
% Тег: #МЕХАНИКА #КИМ_22 #ВАРИАНТ_05
\begin{taskbasic}[code={\label{task:momentum_v5_n22}}]
\textbf{Задача 5.} Тележка массой $40$~кг движется со скоростью, равной по модулю $1$~м/с, слева направо по гладкой горизонтальной дороге. Каким станет модуль скорости тележки, если человек массой $60$~кг запрыгнет на тележку со скоростью, равной по модулю $2$~м/с относительно дороги и направленной слева направо?

\kim{22}
\end{taskbasic}
\answer{1{,}6~м/с}

% === ЗАДАЧА 6 ===
% Тег: #МЕХАНИКА #КИМ_22 #ВАРИАНТ_06
\begin{taskbasic}[code={\label{task:momentum_v6_n22}}]
\textbf{Задача 6.} Тележка массой $60$~кг движется со скоростью, равной по модулю $1$~м/с, слева направо по гладкой горизонтальной дороге. Каким станет модуль скорости тележки, если девочка массой $40$~кг запрыгнет на тележку со скоростью, равной по модулю $2$~м/с относительно дороги и направленной слева направо?

\kim{22}
\end{taskbasic}
\answer{1{,}4~м/с}

% === ЗАДАЧА 7 ===
% Тег: #МЕХАНИКА #КИМ_26 #ВАРИАНТ_15
\begin{taskbasic}[code={\label{task:momentum_v15_n26}}]
\textbf{Задача 7.} Пластилиновый шарик в момент $t = 0$ бросают с горизонтальной поверхности Земли с начальной скоростью $v_0 = 12$~м/с под углом $\alpha = 60^{\circ}$ к горизонту. Одновременно с некоторой высоты над поверхностью Земли начинает падать из состояния покоя другой такой же шарик. Шарики абсолютно неупруго сталкиваются в воздухе. Сразу после столкновения скорость шариков направлена горизонтально. Через какое время после столкновения шариков они упадут на Землю? Сопротивлением воздуха пренебречь. Обоснуйте применимость законов, используемых для решения задачи.

\kim{26}
\end{taskbasic}
\answer{0{,}6~с}

% === ЗАДАЧА 8 ===
% Тег: #МЕХАНИКА #КИМ_26 #ВАРИАНТ_16
\begin{taskbasic}[code={\label{task:momentum_v16_n26}}]
\textbf{Задача 8.} Пластилиновый шарик в момент $t = 0$ бросают с горизонтальной поверхности Земли под углом $\alpha$ к горизонту. Одновременно с некоторой высоты над поверхностью Земли начинает падать из состояния покоя другой такой же шарик. Шарики абсолютно неупруго сталкиваются в воздухе. Сразу после столкновения скорость шариков направлена горизонтально. Время от столкновения шариков до их падения на Землю равно $\tau$. С какой начальной скоростью $v_0$ был брошен первый шарик? Сопротивлением воздуха пренебречь. Обоснуйте применимость законов, используемых для решения задачи.

\kim{26}
\end{taskbasic}
\answer{$\tau g \sqrt{3}$}

% === ЗАДАЧА 9 ===
% Тег: #МЕХАНИКА #КИМ_22 #ВАРИАНТ_25
\begin{taskbasic}[code={\label{task:momentum_v25_n22}}]
\textbf{Задача 9.} Два пластилиновых шарика массами $3m$ и $m$, летящие навстречу друг другу с одинаковыми по модулю скоростями, при столкновении слипаются. Каким был модуль скорости каждого из шариков перед столкновением, если сразу после столкновения скорость шариков стала равной $0{,}5$~м/с? Временем взаимодействия шариков пренебречь.

\kim{22}
\end{taskbasic}
\answer{1~м/с}

% === ЗАДАЧА 10 ===
% Тег: #МЕХАНИКА #КИМ_22 #ВАРИАНТ_26
\begin{taskbasic}[code={\label{task:momentum_v26_n22}}]
\textbf{Задача 10.} Два пластилиновых шарика массами $3m$ и $m$, летящие навстречу друг другу с одинаковыми по модулю скоростями, при столкновении слипаются. Каким станет модуль скорости шариков сразу после столкновения, если перед столкновением модуль скорости каждого из шариков был равен $4$~м/с? Временем взаимодействия шариков пренебречь.

\kim{22}
\end{taskbasic}
\answer{2~м/с}
